% ===================================================================
%  Dodatek T2: Zamknięta forma punktu stałego Koidego
%              i faktoryzacja Z₃
%  TGP v1 — Sesja v41 (2026-04-02/03)
%  Wyniki numeryczne: ex117 (12/12 PASS), ex118 (12/12 PASS),
%                     ex129 (14/14 PASS), ex131 (15/15 PASS),
%                     ex133 (15/17 PASS), ex134 (19/19 PASS),
%                     ex135 (13/16 PASS), ex136 (8/16 PASS),
%                     ex137 (6/16 PASS)
% ===================================================================
\section{Zamkni\k{e}ta forma punktu sta\l{}ego Koidego i~symetria
  $\mathbb{Z}_3$}%
\label{app:T2-fp-algebra}

Niniejszy dodatek rozwija algebr\k{e} r\'ownania punktu sta\l{}ego
wyprowadzonego w~dodatku~T (sekcja~\ref{app:T-koide}).
G\l{}\'ownym wynikiem jest \textbf{dok\l{}adna zamkni\k{e}ta forma}
asymptotycznego kroku wie\.zy leptonowej TGP:
\begin{equation}\label{eq:T2-rstar-exact}
  \boxed{r^* \;=\; \frac{23 + 5\sqrt{21}}{2} \;\approx\; 22{,}9564392374,}
\end{equation}
oraz zwi\k{a}zana z~ni\k{a} \textbf{faktoryzacja}
wielomianu~$P(u) = u^4-4u^3-3u^2-4u+1$ na \emph{sektor fizyczny}
i~\emph{sektor $\mathbb{Z}_3$}:
\begin{equation}\label{eq:T2-factor}
  \boxed{P(u) \;=\; \underbrace{(u^2-5u+1)}_{\text{sektor fiz.}}
                \cdot \underbrace{(u^2+u+1)}_{\text{sektor }\mathbb{Z}_3}.}
\end{equation}

\medskip
\noindent\fbox{\parbox{\linewidth}{%
\textbf{Nota kanoniczna (2026-04-07).}
Niniejszy dodatek u\.zywa Formulacji~B ($g_0^e = 1{,}24915$, $g_0^* = g_0^e$).
W~sformu\l{}owaniu kanonicznym: $K(g) = g^4$, $P(g) = (\beta/7)g^7 - (\gamma/8)g^8$,
$g_0^e = 0{,}86941$, brak punktu duchowego.
Algebra Koide'a i~r\'ownania master F1--F12 s\k{a} niezale\.zne od tego wyboru.}}
\medskip


% -------------------------------------------------------------------
\subsection{R\'ownanie punktu sta\l{}ego Koidego}\label{app:T2-fpeq}
% -------------------------------------------------------------------

Przypomnijmy wynik ex117.
Asymptotyczny krok wie\.zy leptonowej
$s = m_{k+1}/m_k$ stabilizuje si\k{e} do sta\l{}ej $r^*$
wtedy i~tylko wtedy, gdy spe\l{}niony jest \emph{warunek punktu sta\l{}ego}:
\begin{equation}\label{eq:T2-fp-condition}
  Q_K(1,\; s,\; s^2) \;=\; \frac{3}{2},
\end{equation}
gdzie $Q_K(a,b,c) \equiv (\sqrt{a}+\sqrt{b}+\sqrt{c})^2/(a+b+c)$
jest wsp\'o\l{}czynnikiem Koide'go (def.~\ref{def:T-QK}).

Podstawiaj\k{a}c $r_{21}=s$, $r_{31}=s^2$ do definicji~\eqref{eq:T-QK-norm}
i~oznaczaj\k{a}c $u = \sqrt{s}$ (tak \.ze $\sqrt{r_{21}}=u$,
$\sqrt{r_{31}}=u^2$):
\begin{equation}\label{eq:T2-qk-expand}
  Q_K(1, u^2, u^4)
    \;=\; \frac{(1 + u + u^2)^2}{1 + u^2 + u^4}
    \;=\; \frac{3}{2}.
\end{equation}

\begin{proposition}[Wielomian punktu sta\l{}ego]\label{prop:T2-poly}
Warunek~\eqref{eq:T2-fp-condition} jest r\'ownowa\.zny r\'ownaniu
algebraicznemu stopnia czwartego:
\begin{equation}\label{eq:T2-poly}
  P(u) \;\equiv\; u^4 - 4u^3 - 3u^2 - 4u + 1 \;=\; 0,
  \qquad u = \sqrt{r^*} > 0.
\end{equation}

\textbf{Dow\'od}.
Ze wzoru~\eqref{eq:T2-qk-expand} mno\.zymy obie strony przez
$2(1+u^2+u^4)$:
\begin{equation}
  2(1+u+u^2)^2 \;=\; 3(1+u^2+u^4).
\end{equation}
Rozwijaj\k{a}c lewą stron\k{e}:
\begin{equation}
  2\bigl(1 + 2u + 3u^2 + 2u^3 + u^4\bigr)
  \;=\; 3 + 3u^2 + 3u^4.
\end{equation}
Po uproszczeniu:
\begin{equation}
  2 + 4u + 6u^2 + 4u^3 + 2u^4
  - 3 - 3u^2 - 3u^4 \;=\; 0,
\end{equation}
czyli:
\begin{equation}
  -u^4 + 4u^3 + 3u^2 + 4u - 1 \;=\; 0.
\end{equation}
Mno\.z\k{a}c przez $-1$ otrzymujemy r\'ownanie~\eqref{eq:T2-poly}.
\hfill$\square$
\end{proposition}

\begin{remark}[Weryfikacja numeryczna]\label{rem:T2-verify}
R\'ownanie~\eqref{eq:T2-poly} zosta\l{}o zweryfikowane w~\texttt{ex117}
(testy F2 i~F3, 12/12~PASS) oraz \texttt{ex118} (testy G1--G3).
Dla $u^* = \sqrt{r^*} \approx 4{,}79129$:
\[
  |P(u^*)| \;=\; 3{,}55 \times 10^{-14}
  \qquad (\text{precyzja maszynowa}).
\]
\end{remark}


% -------------------------------------------------------------------
\subsection{Palindromiczno\'s\'c i~redukcja}\label{app:T2-palindrome}
% -------------------------------------------------------------------

\begin{proposition}[Wielomian $P$ jest palindromiczny]\label{prop:T2-palindrome}
Wsp\'o\l{}czynniki $P(u) = u^4 - 4u^3 - 3u^2 - 4u + 1$ spe\l{}niaj\k{a}
warunek palindromiczno\'sci:
\begin{equation}
  [a_4, a_3, a_2, a_1, a_0] \;=\; [1,\,-4,\,-3,\,-4,\,1],
  \qquad a_k = a_{4-k},
\end{equation}
tzn.\ $P(u) = u^4 P(1/u)$.

\textbf{Dow\'od}.
\[
  u^4 P\!\left(\tfrac{1}{u}\right)
    = u^4\!\left(\tfrac{1}{u^4} - \tfrac{4}{u^3} - \tfrac{3}{u^2}
      - \tfrac{4}{u} + 1\right)
    = 1 - 4u - 3u^2 - 4u^3 + u^4
    = P(u). \;\hfill\square
\]
\textbf{Konsekwencja}: je\'sli $u_0$ jest pierwiastkiem $P$, to tak\.ze
$1/u_0$ jest pierwiastkiem $P$.
W~szczeg\'olno\'sci pierwiastki $P$ tworz\k{a} pary
$\{u_0,\, 1/u_0\}$.
\end{proposition}

\begin{theorem}[Palindromiczna redukcja]\label{thm:T2-reduction}
Podstawienie $v = u + 1/u$ redukuje r\'ownanie stopnia 4 do stopnia 2:
\begin{equation}\label{eq:T2-vred}
  P(u) \;=\; 0 \quad\Longleftrightarrow\quad
  v^2 - 4v - 5 \;=\; 0
  \quad (v = u + 1/u,\; u \neq 0),
\end{equation}
a~wielomian $v^2-4v-5$ faktoryzuje si\k{e} nad $\mathbb{Z}$:
\begin{equation}\label{eq:T2-vfactor}
  v^2 - 4v - 5 \;=\; (v-5)(v+1).
\end{equation}

\textbf{Dow\'od}.
Dzielimy $P(u)$ przez $u^2 \neq 0$:
\[
  u^2 - 4u - 3 - \frac{4}{u} + \frac{1}{u^2} \;=\; 0.
\]
Grupuj\k{a}c symetrycznie:
\[
  \left(u^2 + \frac{1}{u^2}\right) - 4\!\left(u + \frac{1}{u}\right) - 3
  \;=\; 0.
\]
Korzystamy z~to\.zsamo\'sci $u^2+1/u^2 = (u+1/u)^2 - 2 = v^2-2$:
\[
  (v^2 - 2) - 4v - 3 \;=\; 0,
  \qquad\text{czyli}\quad
  v^2 - 4v - 5 \;=\; 0.
\]
Faktoryzacja nad $\mathbb{Z}$: wsp\'o\l{}czynniki sumuj\k{a} si\k{e}
$5 \cdot (-1) = -5 = a_0$, $5 + (-1) = 4 = -a_1$, st\k{a}d
$(v-5)(v+1) = 0$.
\hfill$\square$
\end{theorem}

\begin{remark}[Podzia\l{} na sektory]\label{rem:T2-sectors}
R\'ownanie $(v-5)(v+1)=0$ wydziela dwa \emph{sektory}:
\begin{itemize}
  \item \textbf{Sektor fizyczny} ($v=5$):
        $u + 1/u = 5 \;\Rightarrow\; u^2 - 5u + 1 = 0$,
        dwa pierwiastki rzeczywiste $(5\pm\sqrt{21})/2$.
  \item \textbf{Sektor $\mathbb{Z}_3$} ($v=-1$):
        $u + 1/u = -1 \;\Rightarrow\; u^2 + u + 1 = 0$,
        dwa pierwiastki zespolone $e^{\pm 2\pi i/3}$.
\end{itemize}
\end{remark}


% -------------------------------------------------------------------
\subsection{Faktoryzacja wielomianu}\label{app:T2-factorization}
% -------------------------------------------------------------------

\begin{theorem}[Faktoryzacja wielomianu FP]\label{thm:T2-factor}
Wielomian punktu sta\l{}ego~\eqref{eq:T2-poly} rozk\l{}ada si\k{e}
nad $\mathbb{Z}$ na dwa nieredukowalne czynniki kwadratowe:
\begin{equation}\label{eq:T2-factor-main}
  u^4 - 4u^3 - 3u^2 - 4u + 1
  \;=\;
  \underbrace{(u^2 - 5u + 1)}_{\substack{\text{sektor fizyczny}\\\text{pierwiastki rzecz.}}}
  \;\cdot\;
  \underbrace{(u^2 + u + 1)}_{\substack{\text{sektor }\mathbb{Z}_3\\\text{pierwiastki zesp.}}}.
\end{equation}

\textbf{Dow\'od}.
Mno\.zymy wprost:
\[
  (u^2-5u+1)(u^2+u+1)
  = u^4 + u^3 + u^2 - 5u^3 - 5u^2 - 5u + u^2 + u + 1
  = u^4 - 4u^3 - 3u^2 - 4u + 1.
\]
Niezmienno\'s\'c nad $\mathbb{Q}$: czynnik $u^2-5u+1$ ma dyskryminant
$\Delta_1 = 25-4 = 21$ (niekwadratowy w~$\mathbb{Q}$), wi\k{e}c jest
nieredukowalny nad $\mathbb{Q}$.  Czynnik $u^2+u+1 = \Phi_3(u)$
jest 3-cyklotomicznym i~r\'ownie\.z nieredukowalny nad $\mathbb{Q}$.
Obydwa czynniki s\k{a} wzajemnie pierwsze ($\gcd = 1$ w~$\mathbb{Q}[u]$).
\hfill$\square$
\end{theorem}

\begin{remark}[To\.zsamo\'s\'c z~sektorami palindromicznej redukcji]
Z~twierdzenia~\ref{thm:T2-reduction}: r\'ownanie $v=5$ daje
$u^2-5u+1=0$, a~$v=-1$ daje $u^2+u+1=0$ ---
dok\l{}adnie dwa czynniki faktoryzacji~\eqref{eq:T2-factor-main}.
Faktoryzacja jest zatem \emph{bezpo\'sredni\k{a} konsekwencj\k{a}}
palindromicznej struktury wielomianu.
\end{remark}


% -------------------------------------------------------------------
\subsection{Zamkni\k{e}ta forma $r^*$}\label{app:T2-closed-form}
% -------------------------------------------------------------------

\begin{theorem}[Zamkni\k{e}ta forma punktu sta\l{}ego]\label{thm:T2-rstar}
Fizyczny punkt sta\l{}y r\'ownania~\eqref{eq:T2-fp-condition}
jest dany wzorem zamkni\k{e}tym:
\begin{equation}\label{eq:T2-rstar}
  r^* \;=\; \frac{23 + 5\sqrt{21}}{2},
\end{equation}
a~odpowiadaj\k{a}cy mu parametr $u^* = \sqrt{r^*}$ spe\l{}nia:
\begin{equation}\label{eq:T2-ustar}
  u^* \;=\; \frac{5 + \sqrt{21}}{2},
  \qquad
  u^* + \frac{1}{u^*} \;=\; 5,
  \qquad
  u^* \cdot \frac{1}{u^*} \;=\; 1.
\end{equation}

\textbf{Dow\'od}.
Z~sektora fizycznego twierdzenia~\ref{thm:T2-factor} pierwiastki
$u^2-5u+1=0$ to:
\[
  u \;=\; \frac{5 \pm \sqrt{21}}{2}.
\]
Dyskryminant $\Delta_1 = 5^2 - 4 \cdot 1 \cdot 1 = 21$.
Poniewa\.z $r^* = u^{*2}$ musi by\'c dodatnie i~$u^* > 1$ (asymptotyczny
krok wie\.zy jest wi\k{e}kszy od 1), wybieramy ga\l{}\k{a}\'z $+$:
\[
  u^* = \tfrac{5+\sqrt{21}}{2} \approx 4{,}7913 > 1,
  \qquad
  \frac{1}{u^*} = \tfrac{5-\sqrt{21}}{2} \approx 0{,}2087 < 1.
\]
Nast\k{e}pnie:
\[
  r^* = \left(\frac{5+\sqrt{21}}{2}\right)^{\!2}
      = \frac{25 + 10\sqrt{21} + 21}{4}
      = \frac{46 + 10\sqrt{21}}{4}
      = \frac{23 + 5\sqrt{21}}{2}.
\]
W\l{}asno\'sci~\eqref{eq:T2-ustar} wynikaj\k{a} z~formul
Viète'a dla $u^2-5u+1=0$:
suma pierwiastk\'ow $= 5$, iloczyn $= 1$.
\hfill$\square$
\end{theorem}

\begin{corollary}[Dok\l{}adno\'s\'c zamkni\k{e}tej formy]\label{cor:T2-precision}
Wzgl\k{e}dna ró\.znica mi\k{e}dzy warto\'sci\k{a}~\eqref{eq:T2-rstar}
a~wynikiem numerycznym brentq (tolerancja $10^{-14}$):
\begin{equation}
  \left|\,r^*_{\rm exact} \;-\; r^*_{\rm brentq}\,\right|
  \;<\; 4 \times 10^{-15}
  \qquad\text{(precyzja maszynowa IEEE-754)}.
\end{equation}
\textbf{Status}: \statuslabel{Twierdzenie}
(analityczne + weryfikacja \texttt{ex118}, test~G3).
\end{corollary}

\begin{remark}[Skąd pochodzi liczba 21?]\label{rem:T2-21}
Dyskryminant $\Delta_1 = 21 = 3 \times 7$.
Czynnik $3$ odzwierciedla liczb\k{e} generacji leptonowych
i~rz\k{a}d grupy $\mathbb{Z}_3$ (por.\ sekcja~\ref{app:T2-Z3}).
Czynnik $7$ nie ma (na razie) bezpo\'sredniej interpretacji w~TGP.
W~szczeg\'olno\'sci $\sqrt{21}$ \textbf{nie jest} prost\k{a} form\k{a}
ani $\varphi^n$, ani $\pi$, ani $e^n$ ---
najmniejsze odchylenie to $3\varphi \approx 4{,}854$ ($\delta = 5{,}9\%$
od $\sqrt{21} \approx 4{,}583$), por.\ \texttt{ex118}, etap~[5].
\end{remark}


% -------------------------------------------------------------------
\subsection{Sektor $\mathbb{Z}_3$ i~symetria Koidego}\label{app:T2-Z3}
% -------------------------------------------------------------------

\begin{proposition}[Czynnik cyklotomiczny $\Phi_3$]\label{prop:T2-cyclo}
Drugi czynnik faktoryzacji~\eqref{eq:T2-factor-main} to:
\begin{equation}
  u^2 + u + 1 \;=\; \Phi_3(u) \;=\; \frac{u^3 - 1}{u - 1},
\end{equation}
gdzie $\Phi_3$ jest \emph{3-cyklotomicznym wielomianem}.
Jego pierwiastki to prymitywne pierwiastki sze\'scienne jedno\'sci:
\begin{equation}\label{eq:T2-omega}
  \omega \;=\; e^{2\pi i/3} \;=\; -\tfrac{1}{2} + \tfrac{\sqrt{3}}{2}\,i,
  \qquad
  \bar\omega \;=\; e^{-2\pi i/3} \;=\; -\tfrac{1}{2} - \tfrac{\sqrt{3}}{2}\,i,
\end{equation}
spe\l{}niaj\k{a}ce $|\omega| = |\bar\omega| = 1$
oraz $\omega^3 = \bar\omega^3 = 1$, $\omega \neq 1$.

\textbf{Dow\'od}.
Bezpo\'srednie sprawdzenie: $\omega^2+\omega+1 = e^{4\pi i/3}+e^{2\pi i/3}+1 = 0$
(suma geometryczna $\sum_{k=0}^{2} e^{2\pi ik/3} = 0$).
$\Phi_3(u) = (u^3-1)/(u-1)$ wynika z~to\.zsamo\'sci
$u^3-1 = (u-1)(u^2+u+1)$.
\hfill$\square$
\end{proposition}

\begin{remark}[$\mathbb{Z}_3$-symetria formu\l{}y Koidego]\label{rem:T2-Z3-koide}
Formu\l{}a Koide'go~\eqref{eq:T-QK} jest jawnie inwariantna
wzgl\k{e}dem permutacji $(\sqrt{m_1}, \sqrt{m_2}, \sqrt{m_3})$,
co generuje grup\k{e} $S_3 \supset \mathbb{Z}_3$.
Podgrupa cykliczna $\mathbb{Z}_3 = \{1, \omega, \omega^2\}$
dzia\l{}a na $(\sqrt{m_k}) \mapsto (\omega^k \sqrt{m_k})$
i~pozostawia $Q_K$ niezmiennym.

Twierdzenie~\ref{thm:T2-factor} pokazuje, \.ze ta symetria jest
\textbf{algebraicznie wbudowana} w~r\'ownanie punktu sta\l{}ego
wie\.zy leptonowej TGP: prymitywne pierwiastki $\mathbb{Z}_3$
wyst\k{e}puj\k{a} w~$P(u)$ dok\l{}adnie jako czynnik $\Phi_3(u) = u^2+u+1$.
Jest to strukturalnie niezale\.zne od sektora fizycznego $(u^2-5u+1)$,
kt\'ory wyznacza $r^*$.
\end{remark}

\begin{remark}[Sektor $v=-1$ a~skok $r^*$]\label{rem:T2-v-minus1}
Z~palindromicznej redukcji (tw.~\ref{thm:T2-reduction}):
$v = \omega + 1/\omega = e^{2\pi i/3} + e^{-2\pi i/3}
= 2\cos(2\pi/3) = -1$ --- dok\l{}adnie ga\l{}\k{a}\'z $v=-1$.
Natomiast sektor fizyczny $v = u^*+1/u^* = 5$ ma swoj\k{a} w\l{}asn\k{a}
liczb\k{e} $5 \in \mathbb{Z}$, ca\l{}kowit\k{a}.
Warto\'sc ta wyznacza
$r^*$ poprzez \eqref{eq:T2-rstar}: $5$ jest \emph{\'sladem}
macierzy Möbiusa $\bigl(\begin{smallmatrix}5&-1\\1&0\end{smallmatrix}\bigr)$
odpowiadaj\k{a}cej przekszta\l{}ceniu $u \mapsto u+1/u = 5$.
\end{remark}


% -------------------------------------------------------------------
\subsection{Wie\.za leptonowa TGP}\label{app:T2-tower}
% -------------------------------------------------------------------

\begin{proposition}[Asymptotyczna wie\.za leptonowa]\label{prop:T2-tower}
Na mocy twierdzenia~\ref{thm:T2-rstar} wie\.za mas leptonowych TGP
jest asymptotycznie geometryczn\k{a} z~ilorazem $r^*$:
\begin{equation}\label{eq:T2-tower}
  m_k \;\xrightarrow{k\to\infty}\;
  m_\tau \cdot \left(\frac{23+5\sqrt{21}}{2}\right)^{\!\!(k-3)},
  \qquad k = 3, 4, 5, \ldots
\end{equation}
Zbieganie jest eksponencjalnie szybkie:
dla $k \geq 7$ krok $r_{k+1,k} = m_{k+1}/m_k$ odbiega
od $r^*$ o~mniej ni\.z $0{,}1\%$ (weryfikacja \texttt{ex117}, test~F7).

Pierwsze generacje powyżej znanych granic eksperymentalnych:
\begin{center}\small
\begin{tabular}{@{}ccrll@{}}
\toprule
$k$ & Lepton & $m_k$ [GeV] & Status & Krok $r_{k,k-1}$ \\
\midrule
1 & $e$     & 0{,}000511  & PDG  & --- \\
2 & $\mu$   & 0{,}1057    & PDG  & $r_{21}=206{,}77$ \\
3 & $\tau$  & 1{,}777     & PDG  & $r_{32}=16{,}82$ \\
4 & $L_4$   & 43{,}7      & WYKL.\ LEP ($>100{,}8$ GeV) & 24{,}59 \\
5 & $L_5$   & \textbf{989} & HL-LHC ($\sim 1\,{\rm TeV}$) & 22{,}63 \\
6 & $L_6$   & $2{,}3 \times 10^4$ & FCC-hh range & $r^*$ \\
$k\geq7$ & $L_k$ & $m_\tau\cdot(r^*)^{k-3}$ & poza LHC & $r^*$ \\
\bottomrule
\end{tabular}
\end{center}

Masa $L_4 = 43{,}7\,\mathrm{GeV}$ jest \textbf{wykluczona} przez LEP
($m > 100{,}8\,\mathrm{GeV}$), co \textbf{potwierdza} dok\l{}adnie
trzy znane generacje leptonowe.
Masa $L_5 = 989\,\mathrm{GeV}$ jest \textbf{pierwsz\k{a} testowalna}
predykcj\k{a} TGP w~zasięgu HL-LHC.
\end{proposition}


% -------------------------------------------------------------------
\subsection{Diagram struktury algebraicznej}\label{app:T2-diagram}
% -------------------------------------------------------------------

\begin{figure}[h]
\begin{center}
\begin{tikzpicture}[
  node distance=1.5cm,
  every node/.style={font=\small},
  box/.style={draw, rounded corners, fill=gray!10, minimum height=0.8cm,
              text centered, inner sep=4pt},
  boxblue/.style={draw=blue!60, rounded corners, fill=blue!8, minimum height=0.8cm,
              text centered, inner sep=4pt},
  boxred/.style={draw=red!50, rounded corners, fill=red!8, minimum height=0.8cm,
              text centered, inner sep=4pt},
  arrow/.style={->, >=Stealth, thick},
]

\node[box, text width=6.5cm] (fp)
  {$Q_K(1,s,s^2)=\frac{3}{2}$ \hfill\emph{warunek FP}};

\node[box, text width=6.5cm, below=0.7cm of fp] (poly)
  {$P(u) = u^4-4u^3-3u^2-4u+1=0$ \hfill ($u=\sqrt{s}$)};

\node[box, text width=6.5cm, below=0.7cm of poly] (palindrome)
  {$P$ palindromiczny: $P(u)=u^4P(1/u)$};

\node[box, text width=6.5cm, below=0.7cm of palindrome] (reduction)
  {$v=u+1/u$: \quad $v^2-4v-5=(v-5)(v+1)=0$};

\node[boxblue, text width=3cm, below left=1.2cm and -0.2cm of reduction] (phys)
  {$v=5$ \\ $u^2-5u+1=0$ \\ $r^*=\dfrac{23+5\sqrt{21}}{2}$};

\node[boxred, text width=3cm, below right=1.2cm and -0.2cm of reduction] (z3)
  {$v=-1$ \\ $u^2+u+1=0$ \\ $\omega=e^{\pm 2\pi i/3}$};

\node[boxblue, text width=3.2cm, below=1.0cm of phys] (tower)
  {Wie\.za leptonowa \\ $m_k = m_\tau\cdot(r^*)^{k-3}$ \\ $L_5=989\,\mathrm{GeV}$};

\node[boxred, text width=3.2cm, below=1.0cm of z3] (sym)
  {Symetria $\mathbb{Z}_3$ \\ permutacja $\sqrt{m_k}$ \\ $\Phi_3(u)=\frac{u^3-1}{u-1}$};

\draw[arrow] (fp) -- (poly) node[midway,right,font=\scriptsize]
  {def.\ $u=\sqrt{s}$, rozwin.};
\draw[arrow] (poly) -- (palindrome) node[midway,right,font=\scriptsize]
  {$[1,-4,-3,-4,1]$};
\draw[arrow] (palindrome) -- (reduction) node[midway,right,font=\scriptsize]
  {$\div u^2$};
\draw[arrow] (reduction) -- (phys) node[midway,left,font=\scriptsize]
  {fiz.};
\draw[arrow] (reduction) -- (z3) node[midway,right,font=\scriptsize]
  {$\mathbb{Z}_3$};
\draw[arrow] (phys) -- (tower);
\draw[arrow] (z3) -- (sym);

\end{tikzpicture}
\end{center}
\caption{Struktura algebraiczna punktu sta\l{}ego Koidego.
  Warunek $Q_K=3/2$ prowadzi do palindromicznego wielomianu $P(u)$,
  kt\'ory przez redukcj\k{e} palindromiczn\k{a} rozk\l{}ada si\k{e}
  na sektor fizyczny (wieża leptonowa, $r^*$)
  i~sektor $\mathbb{Z}_3$ (symetria Koidego, $\Phi_3$).}
\label{fig:T2-algebra-diagram}
\end{figure}


% -------------------------------------------------------------------
\subsection{Podsumowanie wynik\'ow T2}\label{app:T2-summary}
% -------------------------------------------------------------------

\begin{center}\footnotesize
\begin{tabular}{@{}p{5.5cm}lp{5cm}@{}}
\toprule
\textbf{Wynik} & \textbf{Status} & \textbf{Szczeg\'o\l{}y} \\
\midrule
Wielomian FP: $P(u)=u^4-4u^3-3u^2-4u+1$
  & \statuslabel{Twierdzenie}
  & Prop.~\ref{prop:T2-poly}; $|P(u^*)|<4\times10^{-14}$ \\[3pt]
$P(u)$ jest palindromiczny
  & \statuslabel{Twierdzenie}
  & Prop.~\ref{prop:T2-palindrome}; $P(u)=u^4P(1/u)$ \\[3pt]
Redukcja: $v^2-4v-5=(v-5)(v+1)$
  & \statuslabel{Twierdzenie}
  & Tw.~\ref{thm:T2-reduction}; $v=u+1/u$ \\[3pt]
Faktoryzacja $P=(u^2-5u+1)(u^2+u+1)$
  & \statuslabel{Twierdzenie}
  & Tw.~\ref{thm:T2-factor}; nieredukowalny nad $\mathbb{Q}$ \\[3pt]
\textbf{Zamkni\k{e}ta forma} $r^* = \frac{23+5\sqrt{21}}{2}$
  & \statuslabel{Twierdzenie}
  & Tw.~\ref{thm:T2-rstar}; G3: $|\Delta r^*|<4\times10^{-15}$ \\[3pt]
$u^2+u+1 = \Phi_3(u)$ (3-cyklotomiczny)
  & \statuslabel{Twierdzenie}
  & Prop.~\ref{prop:T2-cyclo}; korzenie $e^{\pm 2\pi i/3}$ \\[3pt]
$\mathbb{Z}_3$ symetria algebraicznie wbudowana
  & \statuslabel{Fakt~an.+num.}
  & Rem.~\ref{rem:T2-Z3-koide}; \texttt{ex118} 12/12~PASS \\[3pt]
Wie\.za: $L_5=989\,\mathrm{GeV}$, $L_6=22{,}8\,\mathrm{TeV}$
  & \statuslabel{Predykcja}
  & Prop.~\ref{prop:T2-tower}; \texttt{ex117} 12/12~PASS \\
\bottomrule
\end{tabular}
\end{center}

\bigskip
\noindent\textbf{G\l{}\'owny wniosek.}
Warunek punktu sta\l{}ego wie\.zy leptonowej TGP $Q_K(1,r^*,(r^*)^2)=3/2$
posiada dok\l{}adn\k{a} zamkni\k{e}t\k{a} form\k{e}:
\[
  r^* \;=\; \frac{23+5\sqrt{21}}{2},
\]
a~odpowiadaj\k{a}cy mu wielomian \emph{palindromicznie} rozk\l{}ada si\k{e}
na czynnik fizyczny (wyznaczaj\k{a}cy $r^*$) i~czynnik cyklotomiczny
$\Phi_3(u)$ (enkoduj\k{a}cy symetri\k{e} $\mathbb{Z}_3$ formu\l{}y
Koide'go).  Rezultat ma status \statuslabel{Twierdzenia analitycznego}
zweryfikowanego numerycznie (\texttt{ex117--ex118}, 24/24~PASS).


% -------------------------------------------------------------------
\subsection{Status otwarty T-OP1: Sk\k{a}d $Q_K = 3/2$ dla fizycznych
  leptonów?}\label{app:T2-TOP1}
% -------------------------------------------------------------------

\paragraph{Hierarchia wyja\'snie\'n.}
Twierdzenia T2.1--T2.5 pokazuj\k{a}, \.ze warunek $Q_K=3/2$
jest \emph{algebraicznie konieczny} wzdlu\.z wie\.zy leptonowej TGP
(Poziom~1 poni\.zej).  Pytanie T-OP1 dotyczy poziomu g\l{}\k{e}bszego:
dlaczego fizyczne leptony $(e,\mu,\tau)$ realizuj\k{a} ten warunek
z~dok\l{}adno\'sci\k{a} $\delta Q_K < 0{,}001$?

\begin{center}
\small
\renewcommand{\arraystretch}{1.25}
\begin{tabular}{@{}lll@{}}
\toprule
\textbf{Poziom} & \textbf{Wynik} & \textbf{Status} \\
\midrule
0 & $Q_K(m_e,m_\mu,m_\tau)=3/2$ (PDG, $\delta<0{,}001$)
  & \statuslabel{Obserwacja} \\
1 & FP wie\.zy Koidego $\Rightarrow r^*=(23+5\sqrt{21})/2$
  & \statuslabel{Twierdzenie} (Tw.~\ref{thm:T2-rstar}) \\
2 & $Q_K=3/2 \Leftrightarrow r_{\text{Br}}=\sqrt{2}
       \Leftrightarrow \mathrm{CV}(\sqrt{m})=1$
  & \statuslabel{Fakt an.} (\texttt{ex119}) \\
3 & $A_{\mathrm{tail}}\propto m^{1/4}$
  $\Rightarrow Q_K(A^4)=1{,}500$
  & \statuslabel{Fakt num.} (\texttt{ex125}, $p=0{,}2500$) \\
4 & Sk\k{a}d $Q_K=3/2$ (nie $1{,}472$ z $\varphi$-drabiny)?
  & \statuslabel{T-OP1 OPEN} \\
\bottomrule
\end{tabular}
\end{center}

\paragraph{Hipoteza fazowa PSH (obalona).}
Rozwa\.zono hipotez\k{e}, \.ze asymptotyczna faza solitonu
$\delta_k = \arctan(-C_k/B_k)$ tworzy post\k{e}p
arytmetyczny $\Delta\delta = 2\pi/3$ (symetria $\mathbb{Z}_3$).
Numeryczna weryfikacja (\texttt{ex125}, 10/14~PASS) daje:

\begin{equation}
  \Delta\delta_{12} = \delta_\mu - \delta_e = 167{,}3^\circ,
  \qquad
  \Delta\delta_{23} = \delta_\tau - \delta_\mu = -62{,}9^\circ.
\end{equation}

Warto\'sci te s\k{a} odleg\l{}e od~$\pm 120^\circ$ (post\k{e}pu
$\mathbb{Z}_3$), zatem \textbf{PSH jest obalone}.

\paragraph{Odkrycie pozytywne (ex125).}
Analiza \texttt{ex125} ujawnia, \.ze dla kalibracji
$g_0^k = \varphi^{k-1} g_0^e$ (drabina $\varphi$) zachodzi
$A_{\mathrm{tail}}(g_0^k) \propto m_k^{1/4}$ z~wyk\l{}adnikiem
$p = 0{,}2500$ (log-log OLS, $R^2 \approx 1$).  Wynika st\k{a}d:
\begin{equation}
  Q_K(A_e^4,\, A_\mu^4,\, A_\tau^4)
  \;=\; Q_K(m_e,\, m_\mu,\, m_\tau)
  \;=\; 1{,}5000.
\end{equation}
Oznacza to, \.ze mechanizm amplitudowy $A_{\mathrm{tail}}^4$
\emph{replikuje} warunek Koidego, lecz tylko przez kalibracj\k{e}
drabiny do mas PDG.  Parametr $\xi^*=2{,}553$ (w~miejsce
$\varphi^2=2{,}618$, ró\.znica $2{,}5\%$) pozostaje niewyja\'sniony.

\paragraph{Bliskie trafienie T-OP4.}
Wynik \texttt{ex125} identyfikuje nowe numeryczne trafienie:
\begin{equation}\label{eq:T2-TOP4}
  \frac{m_\mu/m_e}{r^*}
  \;=\; \frac{206{,}768}{22{,}956}
  \;=\; 9{,}007
  \;\approx\; 9 \;=\; N^2 \quad (N=3),
\end{equation}
z~odchyleniem $0{,}078\%$.  Pytanie T-OP4: czy istnieje
algebraiczna to\.zsamo\'s\'c $m_\mu/m_e = N^2 r^*$ w~ramach TGP?

Analiza \texttt{ex129} dostarcza nast\k{e}puj\k{a}cych wynik\'ow:

\begin{enumerate}
  \item \textbf{To\.zsamo\'s\'c analityczna} $u_*^2 = r^*$.
    Z~definicji $u_* = (5+\sqrt{21})/2$ wynika
    $u_*^2 = (46+10\sqrt{21})/4 = (23+5\sqrt{21})/2 = r^*$.
    W~konsekwencji~\eqref{eq:T2-TOP4} mo\.zna przepisa\'c:
    \[
      \frac{m_\mu/m_e}{r^*} \;\approx\; 9
      \;\;\Longleftrightarrow\;\;
      \frac{m_\mu}{m_e} \;\approx\; 9\,u_*^2
      \;\;\Longleftrightarrow\;\;
      \sqrt{\frac{m_\mu}{m_e}} \;\approx\; 3\,u_*,
    \]
    gdzie $3u_* = 3(5+\sqrt{21})/2 = 14{,}374$ wobec
    $\sqrt{m_\mu/m_e} = 14{,}379$ (odchylenie $0{,}039\%$).

  \item \textbf{Formu\l{}a kandydacka TGP.}
    Je\'sli $m_\mu/m_e = N^2 r^*$ dok\l{}adnie, to
    \begin{equation}\label{eq:T2-TOP4-candidate}
      \frac{m_\mu}{m_e}
      \;=\; N^2 r^*
      \;=\; \frac{207 + 45\sqrt{21}}{2}
      \;=\; 206{,}6080\ldots
    \end{equation}
    wobec PDG $206{,}7683$ (odchylenie $0{,}078\%$).

  \item \textbf{K\k{a}t Brannena przy $r_{21}=N^2 r^*$} (wynik \texttt{ex129}, D7).
    W~parametryzacji Brannena
    $\sqrt{m_k} = M(1 + \sqrt{2}\cos(\theta+2\pi k/3))$
    warunek $r_{21}(\theta) = N^2 r^*$ ma rozwi\k{a}zanie
    $\theta_{9r^*} = 132{,}7314^\circ$,
    kt\'ore le\.zy zaledwie $0{,}0014^\circ$ od k\k{a}ta
    $\theta_{\mathrm{PDG}} = 132{,}7328^\circ$.

  \item \textbf{K\k{a}t Brannena z~amplitud solitonu} (wynik \texttt{ex131}).
    \label{it:T2-theta-tgp}
    Skrypt \texttt{ex131} oblicza $\theta_{\mathrm{TGP}}$
    \emph{bezpo\'srednio} z~amplitud ogon\'ow soliton\'ow
    na drabinie~$\varphi$: podstawiaj\k{a}c masy
    $m_k \propto A_{\mathrm{tail}}(g_0^k)^4$ do
    formu\l{}y Brannena~\eqref{eq:T2-TOP4-candidate},
    uzyskuje si\k{e}
    \begin{equation}\label{eq:T2-theta-tgp}
      \theta_{\mathrm{TGP}} \;=\; 132{,}7324^\circ.
    \end{equation}
    Warto\'s\'c ta le\.zy \emph{pomi\k{e}dzy} wynikami \texttt{ex129} i~PDG:
    \begin{equation}\label{eq:T2-theta-ordering}
      \theta_{9r^*} \;=\; 132{,}7314^\circ
      \;<\; \theta_{\mathrm{TGP}} \;=\; 132{,}7324^\circ
      \;<\; \theta_{\mathrm{PDG}} \;=\; 132{,}7328^\circ,
    \end{equation}
    przy czym odst\k{e}py wynosz\k{a}:
    \[
      |\theta_{\mathrm{TGP}} - \theta_{\mathrm{PDG}}| = 0{,}0004^\circ,
      \qquad
      |\theta_{\mathrm{TGP}} - \theta_{9r^*}| = 0{,}0010^\circ.
    \]
    Wszystkie trzy warto\'sci mieszcz\k{a} si\k{e} w~przedziale
    $0{,}0014^\circ$.
    Fakt, \.ze drabina $\varphi$ skalibrowana do mas PDG
    reprodukuje $\theta$ niemal identyczne z~$\theta_{\mathrm{PDG}}$
    i~jednocze\'snie niemal identyczne z~$\theta_{9r^*}$,
    stanowi silne powi\k{a}zanie mi\k{e}dzy trzema pozornie
    niezale\.znymi wielko\'sciami:
    amplitudami soliton\'ow $(A_e, A_\mu, A_\tau)$,
    parametrem Koidego~$r^*$
    oraz mas\k{a} mione $(m_\mu/m_e)$.
\end{enumerate}

\begin{observation}[Zamkni\k{e}cie T-OP4 numeryczne]\label{obs:T2-top4-closure}
  Wyniki \texttt{ex129} i~\texttt{ex131} (15/15~PASS) pokazuj\k{a}:
  \begin{equation}
    \theta_{9r^*} < \theta_{\mathrm{TGP}} < \theta_{\mathrm{PDG}},
    \quad
    |\Delta\theta| < 0{,}0014^\circ.
  \end{equation}
  TGP \emph{numerycznie} odtwarza $\theta_{\mathrm{PDG}}$
  z~dok\l{}adno\'sci\k{a} $|\theta_{\mathrm{TGP}}-\theta_{\mathrm{PDG}}|=0{,}0004^\circ$
  za po\'srednictwem drabiny $\varphi$ i~mechanizmu $A_{\mathrm{tail}}^4$.
  Gdyby istnia\l{}o analityczne wyprowadzenie
  $\theta_{\mathrm{TGP}} = 132{,}7324^\circ$
  z~r\'owna\l{}ania ODE solitonu,
  formu\l{}a $m_\mu/m_e = (207+45\sqrt{21})/2$
  sta\l{}aby si\k{e} \textbf{dok\l{}adn\k{a} predykcj\k{a} TGP}.
\end{observation}

\begin{observation}[Skan $g_0^e$ i~punkt dok\l{}adnego zamkni\k{e}cia
  T-OP4 (\texttt{ex133})]\label{obs:T2-top4-g0e-scan}
  Skrypt \texttt{ex133} (15/17~PASS) wyznacza numerycznie warto\'s\'c
  $g_0^{e,*}$ tak\k{a}, \.ze $\theta_{\mathrm{TGP}}(g_0^{e,*}) = \theta_{9r^*}
  = 132{,}7314^\circ$ dok\l{}adnie:
  \begin{equation}\label{eq:T2-g0e-star}
    g_0^{e,*} = 1{,}249082.
  \end{equation}
  R\'o\.znica wzgl\k{e}dem obserwowanej warto\'sci $g_0^e = 1{,}24915$
  (punkt stały~$\varphi$, ex106) wynosi:
  \begin{equation}
    \Delta g_0^e \;=\; g_0^{e,*} - g_0^e(\mathrm{obs}) \;=\; -0{,}00006
    \quad (= 0{,}005\%).
  \end{equation}
  Czu\l{}o\'s\'c k\k{a}ta na zmian\k{e} $g_0^e$:
  $\mathrm{d}\theta/\mathrm{d}g_0^e \approx 18{,}76\,{}^\circ/\mathrm{unit}$,
  co daje zmian\k{e} $\Delta\theta \approx 0{,}001^\circ$ na zmian\k{e}
  $\Delta g_0^e = 0{,}00006$.
  Ponadto:
  \begin{equation}\label{eq:T2-5-4}
    g_0^e(\mathrm{obs}) = 1{,}24915 \;\approx\; \tfrac{5}{4}
    \quad \text{(b\l\k{a}d } 0{,}068\%),
    \qquad
    g_0^{e,*} = 1{,}249082 \;\approx\; \tfrac{5}{4}
    \quad \text{(b\l\k{a}d } 0{,}073\%).
  \end{equation}
  Zauwa\.zmy, \.ze $\tfrac{5}{4} = 1 + \tfrac{1}{2\alpha} = 1 + \tfrac{1}{4}$
  (przy $\alpha=2$) --- jest to wyk\l{}adnik okre\'slaj\k{a}cy \textbf{\'scian\k{e}
  duchy}: $g^* = e^{-1/(2\alpha)}$.
  Sprawdzono ponadto: $Q_K(g_0^e = 5/4) = 1{,}5007 \approx 3/2$ (\b{l}\k{a}d
  $0{,}05\%$).
\end{observation}

\begin{openproblem}[T-OP4]\label{op:T2-top4}
  Wyprowadzi\'c analitycznie k\k{a}t Brannena
  $\theta_{\mathrm{TGP}} \approx 132{,}7314^\circ$
  z~warunk\'ow TGP: ODE solitonu TGP
  i~asymptotyku ogona $A_{\mathrm{tail}}$.
  Zamkni\k{e}cie da predykcj\k{e}
  $m_\mu/m_e = (207+45\sqrt{21})/2$ (od\-chylenie $0{,}078\%$).
  Wskazanie: skrypt \texttt{ex133} pokazuje, \.ze wystarczy
  przesun\k{a}\'c $g_0^e$ o~zaledwie $0{,}005\%$ (tj.\ o~$\Delta g_0^e =
  -0{,}00006$) od obserwowanej warto\'sci $1{,}24915$, aby osi\k{a}gn\k{a}\'c
  dok\l{}adne zamkni\k{e}cie.
  Kluczowe pytanie: \textbf{dlaczego punkt sta\l{}y $\varphi$-drabiny TGP
  daje $g_0^e \approx 1 + 1/(2\alpha) = 5/4$ z~dok\l{}adno\'sci\k{a} $0{,}07\%$?}
  Liczba $1 + 1/(2\alpha)$ jest wyk\l{}adnikiem \'sciany duchy
  $g^* = e^{-1/(2\alpha)}$, co sugeruje g\l{}\k{e}boki
  zwi\k{a}zek mi\k{e}dzy $\varphi$-drabink\k{a} a~topologi\k{a} potencja\l{}u
  TGP.
  Skrypt \texttt{ex134} (19/19~PASS) identyfikuje nowe wskazanie:
  faza ogona solitonu przy $g_0^e = 5/4$ jest r\'owna $\delta \approx 89{,}1^\circ
  \approx \pi/2$ (odchylenie $0{,}96\%$), zob.\ obserwacj\k{e}~\ref{obs:T2-phase-pi2}.
\end{openproblem}

\begin{observation}[Faza ogona solitonu TGP przy $g_0^e = 5/4$
  (\texttt{ex134})]\label{obs:T2-phase-pi2}
  Skrypt \texttt{ex134} (19/19~PASS) bada analityczn\k{a} struktur\k{e} ODE
  solitonu TGP przy $g_0^e = 5/4 = 1 + 1/(2\alpha)$.
  Ogon solitonu dopasowujemy jako
  $(g-1)\cdot r = B_c\cos r + C_s\sin r$; faza definiujemy jako
  $\delta = \mathrm{arctan2}(C_s, B_c)$.
  G\l{}\'owne wyniki:
  \begin{enumerate}
    \item \textbf{Faza przy $g_0^e = 5/4$:}
      \begin{equation}\label{eq:T2-phase-5-4}
        \delta\!\left(\tfrac{5}{4}\right) \;=\; 89{,}14^\circ
        \;\approx\; \frac{\pi}{2}
        \qquad \text{(odchylenie $0{,}96\%$).}
      \end{equation}
      Ogon solitonu elektronu TGP jest wi\k{e}c \emph{prawie czysto
      sinusoidalny}: $(g-1)\cdot r \approx A_{\mathrm{tail}}\sin r$,
      co odpowiada postaci $g(r) \approx 1 + A_{\mathrm{tail}}\,j_0(r)$,
      gdzie $j_0(r) = \sin r/r$ jest sfery\-czn\k{a} funkcj\k{a} Bessela.

    \item \textbf{Faza liniowego ODE (uwaga: korekta w~\texttt{ex135}):}
      Liniowe ODE $h'' - (2/r)h' + h = 0$ posiada jedynie
      \emph{rosn\k{a}ce} rozwi\k{a}zania dla $r\to\infty$
      (por.\ obserwacja~\ref{obs:T2-ex135}); ogon solitonu jest
      zjawiskiem czysto nieliniowym.
      Jednak w~granicy $A_{\mathrm{tail}}\to0$ (tj.\ $g_0^e\to1^+$)
      numerycznie $\delta \to \approx 90^\circ$:
      \begin{equation}\label{eq:T2-delta0}
        \delta_0 \;\equiv\; \lim_{g_0^e\to1^+}\delta(g_0^e)
        \;\approx\; 90^\circ,
      \end{equation}
      a faktyczna korekta nieliniowa przy $g_0^e = 5/4$ jest
      \emph{ma\l{}a}:
      \begin{equation}\label{eq:T2-delta-phase}
        \Delta\delta \;=\;
        \delta\!\left(\tfrac{5}{4}\right) - 90^\circ
        \;=\; -0{,}864^\circ
        \qquad\text{(nie $\approx\pi/2$, zob.\ \texttt{ex135}).}
      \end{equation}

    \item \textbf{Nieliniowo\'s\'c konieczna:}
      W przybli\.zeniu liniowym $A_{\mathrm{tail}} \propto (g_0^e - 1)$
      warunek $A_\mu/A_e = (9r^*)^{1/4}$ daje $g_0^e \approx 1{,}284$
      (niezgodne z~$5/4$).
      Dopiero wk\l{}ad nieliniowy ($q_{\mathrm{local}} \approx 1{,}094$)
      prowadzi do $g_0^e \approx 5/4$.

    \item \textbf{Korekta drugiego rz\k{e}du:}
      $g_0^{e,*} = 5/4 - 0{,}015\,\varepsilon^2 + \ldots$, gdzie
      $\varepsilon = 1/(2\alpha)$; wsp\'o\l{}czynnik $0{,}015$
      nie ma dotychczas prostej formy analitycznej.

    \item \textbf{Nowe otwarte pytanie (T-OP4c):}
      Dlaczego korekta nieliniowa $\Delta\delta(5/4) = -0{,}86^\circ$
      jest tak ma\l{}a w~por\'ownaniu z~$\delta_0 \approx 90^\circ$,
      tzn.\ sk\k{a}d wynika $|\Delta\delta| \ll \delta_0$
      przy $g_0^e = 5/4$?
      Zob.\ obserwacja~\ref{obs:T2-ex135}.
  \end{enumerate}
\end{observation}

\begin{observation}[Korekta ex134: faza liniowego ODE i~\texttt{ex135}
  (13/16~PASS)]\label{obs:T2-ex135}
  Skrypt \texttt{ex135} (13/16~PASS) koryguje i~uzupe\l{}nia
  analiz\k{e} z~\texttt{ex134}.

  \paragraph{G\l{}\'owna korekta.}
  Liniowe ODE $h'' - (2/r)h' + h = 0$, b\k{e}d\k{a}ce linearyzacj\k{a}
  ODE solitonu TGP przy $g=1$, nie posiada zanikaj\k{a}cych rozwi\k{a}za\'n
  dla $r\to\infty$.  Og\'olne rozwi\k{a}zanie zachowuje si\k{e} jak
  $h(r) \sim r\bigl(C\sin r + D\cos r\bigr)$ (rosn\k{a}ce liniowo).
  Ogon solitonu $(g-1) \sim A\sin(r+\delta)/r$ jest zatem
  zjawiskiem \emph{czysto nieliniowym}: wynika z~globalnej
  struktury ODE, nie z~jego linearyzacji.

  \paragraph{Faza w~granicy ma\l{}ej amplitudy.}
  Numeryczny skan $\delta(g_0^e)$ dla $g_0^e \to 1^+$ (por.\ sek.~[4]
  skryptu) daje
  \begin{equation}\label{eq:T2-delta0-num}
    \delta_0 \;=\; \lim_{A_{\mathrm{tail}}\to0}\delta
    \;\approx\; 90^\circ
    \qquad\bigl(\text{np.\ przy }g_0^e=1{,}01{:}\ \delta=90{,}43^\circ\bigr).
  \end{equation}

  \paragraph{Dok\l{}adne $g_0^e(\delta=90^\circ)$ via brentq.}
  Faza ogona r\'owna si\k{e} dok\l{}adnie $90^\circ$ przy
  \begin{equation}\label{eq:T2-g0e-90deg}
    g_0^e(\delta=90^\circ) \;=\; 1{,}23037
    \qquad\bigl(\text{brentq, } \delta=90{,}0000^\circ \text{ potwierdzone}\bigr),
  \end{equation}
  co r\'o\.zni si\k{e} od $5/4$ o~$-1{,}57\%$.
  Najlepszy kandydat algebraiczny:
  $1 + \varepsilon(1-\varepsilon/3) = 1{,}22917$ (b\l\k{a}d $0{,}098\%$),
  gdzie $\varepsilon = 1/(2\alpha)$.
  Trzy warto\'sci tworz\k{a} hierarchi\k{e}:
  \begin{equation}\label{eq:T2-three-g0e}
    g_0^e(\delta=90^\circ) \;=\; 1{,}23037
    \;<\; \tfrac{5}{4} \;=\; 1{,}25000
    \;\approx\; g_0^{e,*} \;=\; 1{,}24908.
  \end{equation}

  \paragraph{Kwadratowo\'s\'c korekty fazowej.}
  Korekta nieliniowa $\Delta\delta = \delta - 90^\circ$ jest
  kwadratowa w~amplitudzie ogona:
  \begin{equation}\label{eq:T2-Ddelta-quad}
    \Delta\delta(A) \;\approx\;
    -108{,}4^\circ \cdot A^2 + 28{,}0^\circ \cdot A + 0{,}6^\circ,
    \qquad r^2 = 0{,}9997.
  \end{equation}
  Zero tej paraboli: $A \approx 0{,}277$ (co odpowiada
  $g_0^e \approx 1{,}230$~---~zgodnie z~\eqref{eq:T2-g0e-90deg}).
  Maksimum: $\Delta\delta_{\max} \approx 2{,}43^\circ$ przy
  $A \approx 0{,}129$ ($g_0^e \approx 1{,}110$).
  Przy $g_0^e = 5/4$ ($A = 0{,}300$):
  $\Delta\delta(5/4) = -0{,}864^\circ$ --- istotnie \emph{ma\l{}e}.

  \paragraph{Otwarte pytanie T-OP4c.}
  Dlaczego korekta $|\Delta\delta(5/4)| \approx 0{,}86^\circ \ll 90^\circ
  \approx \delta_0$ jest tak ma\l{}a w\l{}a\'snie przy $g_0^e = 5/4
  = 1+1/(2\alpha)$?  Jednocze\'snie $g_0^e = 5/4$ nie jest punktem
  zerowym $\Delta\delta$ (to $g_0^e(\delta=90^\circ) = 1{,}23037$),
  ani punktem sta\l{}ym warunku $\theta_{\mathrm{TGP}} = \theta(9r^*)$
  (to $g_0^{e,*} = 1{,}24908$).
\end{observation}

\begin{observation}[Korekta drugiego rz\k{e}du $c_2 \approx -1/68$
  i~skan~$\alpha$ (\texttt{ex136}, 8/16~PASS)]\label{obs:T2-ex136}
  Skrypt \texttt{ex136} (8/16~PASS) bada rozszerzone zachowanie
  $g_0^{e,*}(\alpha)$ i~wyznacza wsp\'o\l{}czynnik korekty drugiego
  rz\k{e}du przy $\alpha=2$.

  \paragraph{Wysokoprecyzyjne wyznaczenie $g_0^{e,*}$.}
  Dla fizycznego $\alpha = 2$ metoda brentq (10 cyfr) daje
  \begin{equation}\label{eq:T2-g0estar-hprec}
    g_0^{e,*} \;=\; 1{,}249\,081\,6636,
    \qquad \theta_{\mathrm{TGP}}=132{,}731\,439\,00^\circ.
  \end{equation}

  \paragraph{Kandydat na wsp\'o\l{}czynnik korekty.}
  Rozk\l{}adaj\k{a}c $g_0^{e,*} = 1 + \varepsilon + c_2\varepsilon^2
  + O(\varepsilon^3)$ ($\varepsilon = 1/(2\alpha) = 1/4$):
  \begin{equation}\label{eq:T2-c2-value}
    c_2 \;=\;
    \frac{g_0^{e,*} - 1 - \varepsilon}{\varepsilon^2}
    \;=\; -0{,}014693.
  \end{equation}
  Najlepszy kandydat algebraiczny:
  \begin{equation}\label{eq:T2-c2-candidate}
    \boxed{c_2 \;\approx\; -\frac{1}{68} \;=\; -0{,}014706}
    \qquad\bigl(\text{b\l\k{a}d}\;0{,}09\%\bigr),
  \end{equation}
  co mo\.zna zapisa\'c jako $c_2 \approx -\varepsilon/17$ (gdy\.z
  $1/68 = (1/4)/17$).
  Wynikowe przybli\.zenie:
  \begin{equation}\label{eq:T2-g0e-expansion}
    g_0^{e,*} \;\approx\; 1 + \frac{1}{2\alpha}
               - \frac{1}{68\cdot(2\alpha)^2}
               + O\!\left(\frac{1}{\alpha^3}\right).
  \end{equation}

  \paragraph{Zbli\.zenie $g_0^{e,*} \approx 1+1/(2\alpha)$ jest specyficzne
    dla $\alpha = 2$.}
  Skan parametryczny po $\alpha \in \{1{,}5;\,2{,}5;\,\ldots;\,10\}$
  (z~ustalonymi fizycznymi proporcjami $g_0^\mu/g_0^e$, $g_0^\tau/g_0^e$)
  ujawnia, \.ze tylko dla $\alpha=2$ zachodzi zbli\.zenie $g_0^{e,*}
  \approx 1+1/(2\alpha)$ na poziomie $0{,}073\%$.  Dla innych $\alpha$
  odchylenie si\k{e}ga kilku procent lub wi\k{e}cej --- wskazuje to
  na zwi\k{a}zek zbli\.zenia z~\emph{fizyczn\k{a}} warto\'sci\k{a}
  $\alpha=2$ i~sta\l\k{a} Koidego $r^*$.

  \paragraph{Otwarte pytanie T-OP4d.}
  Sk\k{a}d wynika liczba~$17$ w~$c_2 \approx -1/68 = -\varepsilon/17$?
  Czy mo\.zna wyprowadzi\'c $c_2$ analitycznie z~ODE solitonu TGP?
\end{observation}

% -------------------------------------------------------------------
\begin{observation}[Matematyczna $\varphi$-drabina daje
  $g_0^{e,*}(\alpha=2)=1{,}508 \neq 5/4$ (\texttt{ex137}, 6/16~PASS)]
  \label{obs:T2-ex137}
  Skrypt \texttt{ex137} (6/16~PASS) testuje ,,matematyczn\k{a}
  $\varphi$-drabin\k{e}''~--- wariant, w~kt\'orym
  \begin{equation}\label{eq:T2-phi-ladder}
    g_0^\mu \;=\; \varphi\,g_0^e, \qquad
    g_0^\tau \;=\; \varphi^2\,g_0^e,
    \qquad \varphi = \tfrac{1+\sqrt{5}}{2} \approx 1{,}618,
  \end{equation}
  zamiast fizycznych proporcji $R_\mu^{\rm phys}=1{,}6181$,
  $R_\tau^{\rm phys}=2{,}5530$ (ex136).

  \paragraph{G\l{}\'owny wynik: $\varphi^2 \neq R_\tau^{\rm phys}$.}
  Kluczowa r\'o\.znica dotyczy sk\l{}adowej $\tau$:
  \begin{equation}\label{eq:T2-phi2-vs-Rtau}
    \varphi^2 \;=\; 2{,}6180 \;\neq\; R_\tau^{\rm phys} \;=\; 2{,}5530,
    \qquad \Delta R_\tau \;=\; +0{,}065 \;(\approx 2{,}5\%).
  \end{equation}
  Ta pozornie ma\l{}a r\'o\.znica prowadzi do drastycznej zmiany
  punktu sta\l{}ego:
  \begin{align}\label{eq:T2-phi-vs-phys}
    g_0^{e,*}(\varphi\text{-drabina},\;\alpha=2) &= 1{,}508, \\
    g_0^{e,*}(\text{fiz.},\;\alpha=2)            &= 1{,}249,
  \end{align}
  co daje odchylenie $\Delta g_0^{e,*} = +0{,}259$ ($20{,}7\%$).
  Czu\l{}o\'s\'c:
  \begin{equation}\label{eq:T2-sensitivity}
    \frac{\Delta g_0^{e,*}}{\Delta R_\tau}
    \;\approx\; \frac{0{,}259}{0{,}065} \;\approx\; 4{,}0,
  \end{equation}
  czyli $2{,}5\%$ zmiany $R_\tau$ powoduje $20\%$ zmian\k{e}
  $g_0^{e,*}$.

  \paragraph{Zbli\.zenie $g_0^{e,*}\approx 1+1/(2\alpha)$ znika
    pod $\varphi^2$-drabin\k{a}.}
  Dla matematycznej $\varphi$-drabiny \.zadne $\alpha$ ze skanowanego
  zbioru $\{1{,}5;\,2{,}0;\ldots;\,10{,}0\}$ nie spe\l{}nia
  $g_0^{e,*} \approx 1+1/(2\alpha)$ lepiej ni\.z
  $\approx 3\%$ (a dla $\alpha=2$ b\l\k{a}d wynosi $20\%$).
  Korekta $c_2(\alpha=2)=+4{,}13$ jest du\.za i dodatnia,
  w~odro\.znieniu od fizycznego $c_2=-0{,}0147\approx -1/68$.

  \paragraph{Wniosek.}
  Zbli\.zenie $g_0^{e,*}\approx 5/4 = 1+1/(2\alpha)$ jest
  w\l{}asno\'sci\k{a} \emph{fizycznej} proporcji $R_\tau^{\rm phys}
  = \xi^* \approx 2{,}553$, a~nie matematycznej $\varphi$-drabiny.
  Lepton $\tau$~\emph{nie} siedzi w~punkcie $g_0^\tau = \varphi^2
  g_0^e$~--- odchylenie wynosi $2{,}5\%$.

  \paragraph{Otwarte pytanie T-OP4e.}
  Sk\k{a}d pochodzi fizyczna proporcja $R_\tau^{\rm phys}=\xi^*
  \approx 2{,}553$?  Jak wyznaczy\'c $R_\tau(\alpha)$ z~pierwszych
  zasad (punkt sta\l{}y $\varphi$-drabiny z~ex106)?
\end{observation}
